\documentclass[aspectratio=169,t,xcolor=table]{beamer}
\usepackage{amsmath,amssymb}

% ==============================================================================
% STARTER TEMPLATE - CEIA BEAMER DESIGN SYSTEM v3.4
% ==============================================================================
% Use este arquivo para iniciar sua apresentacao de forma rapida e limpa.
% Para compilar:
%   Windows: .\compile.ps1 starter.tex
%   Linux/macOS: ./compile.sh starter.tex
%   Tectonic CLI: tectonic starter.tex
% ==============================================================================

\usetheme{Ceia}

% Metadados Oficiais (Altere com os dados do seu projeto)
\title{TITULO DA SUA APRESENTACAO}
\subtitle{SUBTITULO TECNICO OU INSTITUCIONAL DO PROJETO}
\author{Nome do Autor ou Equipe}
\institute{Centro de Excelencia em Inteligencia Artificial -- UFG}
\date{Goiania, 2026}

\begin{document}

% ==============================================================================
% 1. CAPA OFICIAL (Fundo azul institucional com gotas parametricas TikZ)
% ==============================================================================
\setLayout{cover}
\begin{frame}[noframenumbering]
    \vspace{0.14\paperheight}
    \begin{columns}
        % Coluna Esquerda: Reservada para o padrao de gotas TikZ (Manual CEIA pag. 23)
        \column{0.44\textwidth}
        \phantom{Padrao de Gotas TikZ}
        
        % Coluna Direita: Titulo, Subtitulo, Logo e Autoria
        \column{0.52\textwidth}
        {\fontsize{20}{24}\bfseries\color{white}\raggedright\inserttitle}\par\vspace{0.6em}
        {\fontsize{9.5}{13}\mdseries\color{CEIAGreen}\raggedright\insertsubtitle}\par\vspace{1.2em}
        
        \ceialogowhite[height=1.3cm]\par\vspace{1.2em}
        
        {\fontsize{8.5}{11}\mdseries\color{white!85}\insertauthor}\par\vspace{0.2em}
        {\fontsize{7.5}{9.5}\mdseries\color{white!65}\insertinstitute}
    \end{columns}
\end{frame}

% ==============================================================================
% 2. CONTEXTO E OBJETIVOS (Layout Content - Cards Modulares)
% ==============================================================================
\setLayout{content}
\begin{frame}{CONTEXTO E MOTIVACAO}{PILARES ESTRATEGICOS DO PROJETO}
    \vspace{0.5em}
    \begin{ceiacards}[3]
        \ceiacarditem[1]{Pesquisa Aplicada}{Modelos de estado da arte treinados com rigor cientifico e reproducibilidade.}
        \ceiacarditem[2]{Soberania de Dados}{Execucao e orquestracao on-premise com total aderencia a LGPD e seguranca.}
        \ceiacarditem[3]{Engenharia e Escala}{Pipelines resilientes, tipagem estrita de contratos e monitoramento em tempo real.}
    \end{ceiacards}
\end{frame}

% ==============================================================================
% 3. ARQUITETURA MULTI-AGENTE (DSL Declarativa TikZ sem coordenadas)
% ==============================================================================
\begin{frame}{ARQUITETURA DA SOLUCAO}{PIPELINE DECLARATIVO COM GATES DE VALIDACAO}
    \vspace{0.5em}
    % Sintaxe: \ceiapipelinethree{Label 1}{role}{Label 2}{role}{Label 3}{role}
    % Roles disponiveis: agent, gate, tool, db
    \ceiapipelinethree{Elicitacao \&\\Parser Zod}{agent}{[G1] Gate de\\Homologacao}{gate}{Execucao \&\\Observabilidade}{tool}
    
    \vspace{1.2em}
    \ceiabanner[ok]{Contratos Formais Garantidos}{A validacao deterministica elimina 100\% dos incidentes de parsing silencioso em runtime.}
\end{frame}

% ==============================================================================
% 4. TABELA EXECUTIVA E RESULTADOS (Tabela com marcadores vetoriais)
% ==============================================================================
\begin{frame}{BENCHMARKS E IMPACTO}{AVALIACAO COMPARATIVA DE DESEMPENHO}
    \vspace{0.5em}
    \begin{ceiamatrixtable}{l c c c}
        \textbf{Requisito Tecnico} & \textbf{Abordagem Legada} & \textbf{Cloud Generica} & \textbf{Padrao CEIA} \\
        \midrule
        \ceiaicon{shield}\hspace{4pt} Governanca \& LGPD & \matrixfail & \matrixwarn & \matrixok (Soberana) \\
        \ceiaicon{database}\hspace{4pt} Rastreabilidade OpenTelemetry & \matrixfail & \matrixok & \matrixok (100\%) \\
        \ceiaicon{clock}\hspace{4pt} Latencia p95 em Producao & 3.4s & 1.2s & \bestmetric{110ms} \\
    \end{ceiamatrixtable}
    
    \vspace{1.0em}
    \begin{columns}[T]
        \begin{column}{0.48\textwidth}
            \ceiastatfour{99.8\%}{Disponibilidade}{3.2x}{Throughput}{0}{Vazamentos}{110ms}{Latencia p95}
        \end{column}
        \begin{column}{0.48\textwidth}
            \ceiabanner[info]{Impacto Direto}{Reducao de 65\% nos custos operacionais com throughput 3 vezes superior.}
        \end{column}
    \end{columns}
\end{frame}

% ==============================================================================
% 5. DIVISOR DE CAPITULO (Layout Section paramétrico oficial)
% ==============================================================================
% Sintaxe: \ceiasectionframe[right|bottom]{Numero}{TITULO}{Subtitulo}
\ceiasectionframe{02}{SINTESE E PROXIMOS PASSOS}{Consolidacao dos Resultados e Cronograma de Entrega}

% ==============================================================================
% 6. CONCLUSAO EXECUTIVA (Card Monumental Anti-Vacuo e Chips de Alto Contraste)
% ==============================================================================
\setLayout{content}
\begin{frame}{VEREDITO EXECUTIVO}{SINTESE DOS RESULTADOS HOMOLOGADOS}
    \vspace{0.4em}
    
    \begin{ceiasummarypunch}{VEREDITO DE ENGENHARIA \& PESQUISA}{Arquitetura Homologada para Operacao Critica em Producao}
        A migracao para orquestracao declarativa e contratos formais garantiu estabilidade deterministica, eliminando riscos de regressao e assegurando total conformidade regulatoria.
        \par\vspace{12pt}
        \ceiachip[CEIAGreen]{shield}{Conformidade LGPD \& EU AI Act}\hfill
        \ceiachip[CEIABrightBlue]{server}{Soberania On-Premise}\hfill
        \ceiachip[CEIADarkBlue]{clock}{Latencia p95 $<$ 120ms}
    \end{ceiasummarypunch}
\end{frame}

% ==============================================================================
% 7. ENCERRAMENTO OFICIAL (Layout closing)
% ==============================================================================
\setLayout{closing}
\begin{frame}[noframenumbering]
    \begin{tikzpicture}[remember picture, overlay]
        % Centro: OBRIGADO :)
        \node[anchor=center, align=center] at ($(current page.center)+(0, 0.02\paperheight)$) {
            {\fontsize{44}{50}\selectfont\bfseries\color{white} OBRIGADO :)}
        };

        % Rodape Institucional
        \node[anchor=south, align=center] at ($(current page.south)+(0, 0.75)$) {
            {\fontsize{11}{14}\selectfont\color{CEIAGreen}\bfseries ceia.ufg.br \quad\textbar\quad contato@ceia.ufg.br}\\[6pt]
            {\fontsize{8}{10}\selectfont\color{white!70} CENTRO DE EXCELENCIA EM INTELIGENCIA ARTIFICIAL --- UFG}
        };
    \end{tikzpicture}
\end{frame}

\end{document}
